\lecture{5}{26. September 2025}{Partial Differential Equations}

\begin{problem}{5.1}
  Write down the general expression for a linear first order PDE for $u(x,y,t)$.
\end{problem}

\begin{derivation}
  From the Lecture (p. 1) the general expression for a linear first order PDE for $u(x,y)$ is
  \begin{equation*}
    a_0(x,t) + a_1(x,t)u(x,t) + a_2(x,t) \frac{\partial u(x,t)}{\partial x} + a_3(x,t) \frac{\partial u(x,t)}{\partial t} = 0
  .\end{equation*}
  Extending this pattern we can write the general expression for a linear first order ODE for $u(x,y,t)$ as
\end{derivation}

\finalresult{
    a_0(x,y,t) + a_1(x,y,t)u(x,y,t) + a_2(x,y,t) u_x(x,y,t) + a_3(x,y,t) u_y(x,y,t) + a_4(x,y,t) u_t(x,y,t) = 0
}

\begin{problem}{5.2}
  Write down the general expression for a linear second order PDE for $u(x,y)$.
You may assume that $u(x,y)$ is twice continuously differentiable.
\end{problem}

\begin{derivation}
  Extending the logic from the previous answer, the general expression for a linear second order PDE for $u(x,y)$ is:
\end{derivation}

\finalresult{
\begin{aligned}
  a_0(x,y) &+ a_1(x,y)u(x,y) + a_2(x,y) u_x(x,y) + a_3 (x,y) u_y(x,y) \\
                            & \qquad + a_4(x,y) u_{x x}(x,y) + a_5(x,y) u_{yy}(x,y) + a_6(x,y) u_{xy}(x,y) = 0
\end{aligned}
}

\begin{problem}{5.3}
  Consider the general homogeneous first order PDE for $u(x,t)$ with solutions $u_1(x,t)$ and $u_2(x,t)$.
Show that $u(x,t) = c_1 u_1 (x,t) + c_2 u_2 (x,t)$ is also a solution for arbitrary constants $c_1$ and $c_2$.
\end{problem}

\begin{derivation}
  From the Lecture (p. 2) the general homogeneous first order PDE for $u(x,t)$ is
  \begin{equation*}
    a_1(x,t) u(x,t) + a_2(x,t) u_x(x,t) + a_3(x,t) u_t(x,t) = 0
  .\end{equation*}
  We will see if this holds for arbitrary constants $c_1$ and $c_2$ as
  \begin{align*}
    & \quad a_1 (c_1 u_1 + c_2 u_2) + a_2 \frac{\partial c_1 u_1 + c_2 u_2}{\partial x} + a_3 \frac{\partial c_1 u_1 + c_2 u_2}{\partial t} \\
    &= a_1 c_1 u_1 + a_1 c_2 u_2 + a_2 \frac{\partial c_1 u_1}{\partial x} + a_2 \frac{\partial c_2 u_2}{\partial x} + a_3 \frac{\partial c_1 u_1}{\partial t} + a_3 \frac{\partial c_2 u_2}{\partial t}  \\
    &= c_1 \left( a_1 u_1 + a_2 \frac{\partial u_1}{\partial x} + a_3 \frac{\partial u_1}{\partial t}  \right) + c_2 \left( a_1 u_2 + a_2 \frac{\partial u_2}{\partial x} + a_2 \frac{\partial u_2}{\partial t}  \right) \\
    &= c_1 \left( 0 \right) + c_2 \left( 0 \right) \\
    &= 0
  .\end{align*}
  Hence $u(x,t) = c_1 u_1(x,t) + c_2 u_2(x,t)$ is also a solution of the general homogeneous first order PDE for $u(x,t)$.
\end{derivation}

\newpage

\begin{problem}{5.4}
  Find the solution of the one-dimensional wave equation
\[ 
u_{tt} (x,t) = u_{x x} (x,t), \quad 0 \leq x \leq 1
\]
subject to the boundary conditions
\[ 
u(0,t) = u(1,t) = 0, \quad \text{for all } t
\]
and the initial condition
\[ 
u(x,0) = k_0 \sin \left( 3 \pi x \right), \quad u_t(x,0) = 0, \quad 0 \leq x\leq 1,
\]
where $k_0$ is a constant.
You may start your calculations directly from the general solution given in the Lecture.
\end{problem}

\begin{derivation}
  From the Lecture (p. 8) the general solution to the one-dimensional wave equation is
  \begin{equation*}
    u(x,t) = \sum_{n = 1}^{\infty} \left( B_n \cos \left( \lambda_n t \right) + B_n^{*} \sin \left( \lambda_n t \right) \right) \sin \left( n \pi x \right) \quad \lambda = n \pi
  .\end{equation*}
  From the initial elongation (p. 9) we get
  \begin{align*}
    u(x,0) &= \sum_{n = 1}^{\infty} B_n \sin \left( n \pi x \right) \\
    &= k_0 \sin \left( 3\pi x \right) \\
    \implies \quad B_3 &= k_0
  .\end{align*}
  And all other coefficients $B_n$ are zero.

  From the initial velocity (p. 10) we get
  \begin{align*}
    u_t(x,0) &= \sum_{n = 1}^{\infty} B_n^{*}\lambda_n \sin \left( n \pi x \right) = 0\\
    \implies \quad B_n^{*} &= 0 \text{ for all } n
  .\end{align*}
  Hence, the particular solution is
\end{derivation}

\finalresult{
    u(x,t) = k_0 \cos (3\pi t) \sin (3 \pi x)
}


\begin{problem}{5.5}
  Find the solution of the one-dimensional wave equation
\[ 
u_{tt} (x,t) = u_{xx}(x,t), 0 \leq x \leq 1
\]
subject to the boundary conditions
\[ 
u(0,t) = u(1,t) = 0, \quad \text{for all } t
\]
and the initial condition
\[ 
u(x,0) = k_0 \sin \left( 3 \pi x \right), \quad u_t (x,0) = 3\pi k_0 \sin \left( 3 \pi x \right), 0 \leq x \leq 1,
\]
where $k_0$ is a constant.
You may start your calculation directly from the general solution given in the Lecture.
\end{problem}

\begin{derivation}
  As the equation, solution and initial elongation are the same as for Problem 5.4 we know that $B_3 = k_0$.
  We therefore start with the initial velocity (which differs in this case).
  We get
  \begin{align*}
    u_t(x,0) &= \sum_{n = 1}^{\infty} B_n^{*}\lambda_n \sin \left( n \pi x \right) \\
    &= 3 \pi k_0 \sin \left( 3 \pi x \right) \\
    \implies \quad B_3^{*} &= k_0
  .\end{align*}
  And all other coefficients $B_n^{*} = 0$.
  Hence,
\end{derivation}

\finalresult{
u(x,t) = k_0 \sin (3\pi x) ( \cos (3\pi t) + \sin(3\pi t))
}
